\chapter{Persistence, Admissibility, and Structural Exclusion}\label{ch:persist}

Does \(G\) always collapse to 1? If it did, regime~I would be the whole story. It does not. This chapter introduces the integer coordinate \(\gamma\) that makes persistent orbits tractable, constructs persistent regimes for infinitely many \(g\), shows that universal collapse is false, and proves that certain moduli can never persist.

\section{\texorpdfstring{The \(\gamma\)-coordinate}{The gamma-coordinate}}

For a realised \(x\) with \(g=G(x)\) and \(L=L(x)\), Proposition~\ref{prop:frame} gives \(x=\gamma L^2/g\) with \(\gamma\in\Z_{>0}\). This coordinate evolves by simple rules.

\begin{theorem}[The \(\gamma\)-framework]\rec{F.T44}\scope{C}\label{thm:gamma}
Let \(x\) be realised, \(g=G(x)\), \(L=L(x)\), \(\gamma=\gamma(x)\).
\begin{enumerate}[label=(\alph*),nosep]
\item \(x=\gamma L^2/g\) with \(\gamma\in\Z_{>0}\), \(b\mid g\), \(\gcd(L,g)=1\).
\item The number \(\gamma L^2/g\) is an integer divisible by 4 iff \(4g\mid\gamma\) (\(L\) odd) or \(g\mid\gamma\) (\(L\) even, \(g\) odd). (For \(g>1\) this cannot occur at a state of the orbit: it would force \(G=1\).)
\item If the step \(x\to x'\) is \(G\)-preserving and \(r=L(x')/L(x)\in\Z_{\ge1}\), then \(\gamma(x')=(\gamma+1)/r^2\). On a plateau \(\gamma'=\gamma+1\).
\item (\emph{Congruence advance}) If \(g\) is even, then \(r\) is odd and \(\gamma'\equiv\gamma+1\ (\mathrm{mod}\ 8)\). If \(3\mid g\), then \(\gamma'\equiv\gamma+1\ (\mathrm{mod}\ 3)\). If \(6\mid g\), then \(\gamma'\equiv\gamma+1\ (\mathrm{mod}\ 24)\).
\item The reduced denominator of \(x\) is \(b=g/\gcd(g,\gamma)\); so \(g\mid\gamma\) iff \(x\) is an integer.
\end{enumerate}
\end{theorem}
\begin{proof}
(a) is Proposition~\ref{prop:frame}. (b) \(\gcd(L,g)=1\). If \(L\) is odd, \(\gcd(L^2,4g)=1\) and \(4g\mid\gamma L^2\) iff \(4g\mid\gamma\). If \(L\) is even then \(g\) is odd (for even \(g\) all \(d\) are odd), \(L^2=4(L/2)^2\) with \(\gcd(L/2,g)=1\), and \(4g\mid4\gamma(L/2)^2\) iff \(g\mid\gamma\). If \(G(x)>1\) the conclusion \(G=1\) of Theorem~\ref{thm:G-arith}(c) is impossible. (c) With \(G'=g\), \(gx'=y+L^2\) and \(\gamma'=gx'/L'^2=(\gamma+1)L^2/L'^2=(\gamma+1)/r^2\). (d) \(\gamma'r^2=\gamma+1\). If \(g\) is even, every modulus is odd, so \(r\) is odd and \(r^2\equiv1\ (\mathrm{mod}\ 8)\). If \(3\mid g\), all moduli are \(\equiv1\ (\mathrm{mod}\ 3)\), so \(3\nmid r\) and \(r^2\equiv1\ (\mathrm{mod}\ 3)\); if \(6\mid g\), \(r\) is prime to 6 and \(r^2\equiv1\ (\mathrm{mod}\ 24)\). (e) \(\gcd(L^2,g)=1\), so the reduced denominator of \(\gamma L^2/g\) is \(g/\gcd(g,\gamma)\).
\end{proof}

Call an orbit \emph{persistent} if \(G_n=g>1\) for all \(n\ge n_0\); since \(G_n\) is non-increasing through divisors, \(G_\infty\) is attained, and every orbit with \(G_\infty>1\) is eventually persistent. In a persistent orbit every growth event is \(G\)-preserving (Theorem~\ref{thm:transition}), so Theorem~\ref{thm:gamma}(c) applies at every step: the orbit is the pair \((\gamma_n,L_n)\), with \(\gamma\) advancing by 1 on plateaus and dropping to \((\gamma+1)/r^2\) at carries.

\section{Persistent regimes}

\begin{theorem}[The regime \(P_g\)]\rec{F.T38}\scope{C}\label{thm:Pg}
Let \(g\ge5\) with \(\gcd(g,6)=1\) and \(s=2^km\) with \(k\ge1\), \(m\) odd, \(\gcd(m,g)=1\). Assume
\begin{itemize}[nosep]
\item[(H1)] the residues modulo \(g\) of the divisors of \(m\) contain the subgroup \(\langle2\rangle\le(\Z/g)^\times\);
\item[(H2)] every prime power \(q^e\,\|\,m\) divides some divisor \(d\) of \(2^km\) with \(d\equiv1\ (\mathrm{mod}\ g)\);
\item[(H3)] \(\gcd\{(d-1)/g:\ d\mid s,\ d\ge2,\ d\equiv1\ (\mathrm{mod}\ g)\}=1\).
\end{itemize}
Put \(x_0=s^2t/g\) with \(t\in\{1,2,3\}\), and write \(s_n=s\cdot(\text{product of the factors }\sigma)\), \(t_{n+1}=\sqf(t_n+1)\) with \(\sigma_n^2\,\|\,t_n+1\) (\(\sigma=2\) exactly when \(t_n=3\), else \(1\)). Then for every \(n\ge0\):
\begin{enumerate}[label=(\roman*),nosep]
\item \(x_n=s_n^2t_n/g\) in lowest terms, \(t_n\in\{1,2,3\}\);
\item \(\Dx(x_n)=\{d\mid s_n:\ d\ge2,\ d\equiv1\ (\mathrm{mod}\ g)\}\), \(L(x_n)=s_n\), \(G(x_n)=g\);
\item \(x_{n+3}=4x_n\), and \(G(x_n)=g\) for all \(n\): the orbit never reaches \(G=1\).
\end{enumerate}
\end{theorem}
\begin{proof}
Induction on \(n\). Suppose (i) at \(n\), with \(2^{k_n}\,\|\,s_n\), \(k_n\ge k\), and odd part \(m\) (unchanged, since \(\sigma\in\{1,2\}\)). Because \(g\) is prime to 6, \(\gcd(s,g)=\gcd(t,g)=1\), so \(x_n=s_n^2t_n/g\) is reduced with \(b=g\). As \(t_n\) is squarefree, the square-root part is \(s_n\) and, by Lemma~\ref{lem:sqrtpart}, \(\Dx(x_n)=\{d\mid s_n,d\ge2,d\equiv1\ (\mathrm{mod}\ g)\}\). \emph{\(L=s_n\):} for the 2-part, (H1) gives \(u\mid m\) with \(u\equiv2^{-k_n}\ (\mathrm{mod}\ g)\), so \(d=2^{k_n}u\in\Dx\); for \(q^e\,\|\,m\), (H2) gives a modulus divisible by \(q^e\) (it divides \(s_n\) because \(s_n\) is a multiple of \(2^km\)). \emph{\(G=g\):} \(g\mid G\) as every \(d\equiv1\ (\mathrm{mod}\ g)\); \(G(x_n)\mid G(x_0)\) by Theorem~\ref{thm:mono}; and \(G(x_0)=g\gcd\{(d-1)/g\}=g\) by (H3). Hence \(\Dcoll=s_n^2/g\) and
\[
x_{n+1}=\frac{s_n^2t_n+s_n^2}{g}=\frac{s_n^2(t_n+1)}{g}=\frac{s_{n+1}^2t_{n+1}}{g},
\]
with \(t_{n+1}\in\{2,3,1\}\) prime to \(g\). So (i) and (ii) hold at \(n+1\). Over three steps the factors \(\sigma\) are \(1,1,2\) in some phase, so \(s\) doubles and \(t\) returns: \(x_{n+3}=4x_n\).
\end{proof}

\begin{proposition}[Existence for every \(g\) prime to 6]\rec{F.T38a}\scope{C}\label{prop:T38a}
Let \(g\ge5\), \(\gcd(g,6)=1\), \(q\) the least odd prime with \(q\equiv2\ (\mathrm{mod}\ g)\) (Dirichlet), \(o=\ord_g(2)\), \(g+1=2^vw\) with \(w\) odd, \(m=wq^o\), \(k=\max(v,1)\). Then (H1)--(H3) hold for \(s=2^km\).
\end{proposition}
\begin{proof}
\(\gcd(w,g)=1\) as \(w\mid g+1\), and \(\gcd(q,g)=\gcd(2,g)=1\). A prime divisor of \(g+1\) is at most \(g+1\), and a prime \(\equiv2\ (\mathrm{mod}\ g)\) at most \(g+1\) is \(2\); so the odd prime \(q\) does not divide \(g+1\), and \(q^o\,\|\,m\). (H1): the divisors \(q^i\equiv2^i\) (\(0\le i<o\)) exhaust \(\langle2\rangle\). (H2): \(q^o\equiv2^o=1\ (\mathrm{mod}\ g)\) is a modulus divisible by \(q^o\); each prime power of \(w\) divides \(g+1\mid s\), and \(g+1\equiv1\). (H3): \(d=g+1\) has \((d-1)/g=1\).
\end{proof}

The Dirichlet step can be avoided altogether.

\begin{theorem}[Elementary construction]\rec{F.T42}\scope{C}\label{thm:T42}
Let \(g\ge5\), \(\gcd(g,6)=1\), \(o=\ord_g(2)\), \(s=(g+1)(g+2)^o\) and \(x_0=s^2/g\). Then \(x_n=s_n^2t_n/g\) with \(t_n\in\{1,2,3\}\), \(G(x_n)=g\) and \(x_{n+3}=4x_n\) for all \(n\).
\end{theorem}
\begin{proof}
Let \(k=\vv_2(g+1)\ge1\) and \(m=s/2^k=\mathrm{odd}(g+1)\,(g+2)^o\), odd and prime to \(g\) (\(g\) is odd and \(\gcd(g+1,g)=\gcd(g+2,g)=1\)). We verify the hypotheses of Theorem~\ref{thm:Pg}. (H1): the divisors \((g+2)^j\) of \(m\) are \(\equiv2^j\ (\mathrm{mod}\ g)\), so their residues contain \(\langle2\rangle\). (H2): since \(\gcd(g+1,g+2)=1\), a prime power \(q^e\,\|\,m\) lies in \(g+1\) (and \(g+1\equiv1\) is a modulus dividing \(s\)) or in \((g+2)^o\), which is itself a modulus: \((g+2)^o\equiv2^o=1\). (H3): \(d=g+1\) gives \((d-1)/g=1\).
\end{proof}

\emph{Instances.} Direct iteration of \(C\) confirms Theorem~\ref{thm:T42} for \(g=5\), \(7\), \(11\), \(13\), \(17\), \(19\), \(23\), \(25\), \(29\), \(35\), \(49\), \(55\), \(65\), \(77\) (twelve steps each, in exact rational arithmetic, together with the hypotheses (H1)--(H3)). For \(g=5\): \(o=4\), \(s=6\cdot7^4=14406\), and \(R(x_0)=\#\{d\mid s:\ d\ge2,\ d\equiv1\ (\mathrm{mod}\ 5)\}\).

\begin{theorem}[Persistent regimes with \(3\mid g\)]\rec{F.T46}\scope{C}\label{thm:Pstar}
Let \(g\ge5\) be odd (\(3\mid g\) allowed), \(s=2^km\) with \(m\) odd, \(\gcd(m,g)=1\), \(o=\ord_g(2)\), and assume (H1)--(H3) of Theorem~\ref{thm:Pg} and
\begin{itemize}[nosep]
\item[(H4)] for \(t\in\{1,2,3\}\) and \(b_t=g/\gcd(g,t)\), every \(d=2^ja\) (\(a\mid m\), \(j\ge0\)) with \(d\equiv1\ (\mathrm{mod}\ b_t)\) satisfies \(d\equiv1\ (\mathrm{mod}\ g)\).
\end{itemize}
Then for \(x_0=s^2/g\) the conclusions (i)--(iii) of Theorem~\ref{thm:Pg} hold, and so does the periodic increment of Theorem~\ref{thm:T43}.
\end{theorem}
\begin{proof}
As for Theorem~\ref{thm:Pg}; the single new point is that \(x_n=s_n^2t_n/g\) has reduced denominator \(b_t=g/\gcd(g,t)\) (because \(\gcd(s,g)=1\) and \(t\) is squarefree), so \(\Dx(x_n)=\{d\mid s_n:\ d\equiv1\ (\mathrm{mod}\ b_t)\}\), which equals \(\{d\mid s_n:\ d\equiv1\ (\mathrm{mod}\ g)\}\) by (H4) (the residue of \(d\) depends only on \(j\bmod o\)). Then \(G(x_n)\mid G(x_0)=g\), \(b\mid G\), and \(g\mid G\) give \(G=g\) as before.
\end{proof}

\emph{Instances.} \(g=15,39,51,57,87\), with \(s=(g+1)(g+2)^o\) (direct iteration, 15 steps each; (H4) checked by exhaustion). So \(G_\infty>1\) occurs with \(\gcd(G_\infty,6)=3\).

\begin{remark}[Why \(3\mid g\) works and \(2\mid g\) does not]
For \(g=15\) and \(t=3\) the denominator is \(b=5\), and the divisors of \(s\) lie in \(\langle2\rangle\bmod15\), on which reduction modulo 5 is injective (\(\ord_5(2)=\ord_{15}(2)=4\)); so no new modulus appears. For even \(g\), every modulus \(d\equiv1\ (\mathrm{mod}\ g)\) is odd, powers of 2 are never moduli, and \(L\ne s\); the \((s,t)\)-description of Theorem~\ref{thm:Pg} breaks down. Even \(g\) is governed instead by the carry structure of Chapter~\ref{ch:gamma}.
\end{remark}

\begin{theorem}[Periodic increment of the fibre]\rec{F.T43}\scope{C}\label{thm:T43}
In any orbit satisfying Theorem~\ref{thm:Pg} or~\ref{thm:Pstar}, let \(s_n=2^{k_n}m\) (\(m\) odd, fixed), \(o=\ord_g(2)\), and \(A_g(m)=\#\{a\mid m:\ a\bmod g\in\langle2\rangle\}\). Then
\[
R(x_n)=\#\{(a,j):a\mid m,\ 0\le j\le k_n,\ 2^ja\equiv1\ (\mathrm{mod}\ g)\}-1,\qquad k_{n+3o}=k_n+o,
\]
and \(R(x_{n+3o})-R(x_n)=A_g(m)\) for every \(n\).
\end{theorem}
\begin{proof}
\(\Dx(x_n)=\{d\mid s_n:d\ge2,\ d\equiv1\ (\mathrm{mod}\ g)\}\), and \(d=2^ja\) uniquely (\(a\mid m\) odd, \(j\le k_n\)); the \(-1\) removes \(d=1\). Every three steps \(s\) is multiplied by exactly 2, so \(k_{n+3o}=k_n+o\). The new pairs have \(j\in(k_n,k_n+o]\). Since \(2^j\bmod g\) has exact period \(o\), over \(o\) consecutive values of \(j\) the residue \(2^j\) meets each element of \(\langle2\rangle\) exactly once; so for each \(a\) with \(a\bmod g\in\langle2\rangle\) exactly one \(j\) gives \(2^ja\equiv1\), and none for other \(a\).
\end{proof}

\begin{theorem}[Reach set and growth in \(P_g\)]\rec{F.T40}\scope{C}\label{thm:T40}
In the orbits of Theorems~\ref{thm:Pg} and~\ref{thm:Pstar} the reach set is
\(U=\{d\ge2:\ d\mid2^km\text{ for some }k,\ d\equiv1\ (\mathrm{mod}\ g)\}\).
It contains no multiple of \(g\) and none of \(2,3,\dots,g\), and \(U(X)=O(\log X)\) (density zero), while \(R(x_n)\) grows linearly with slope \(A_g(m)/(3o)\) per step.
\end{theorem}
\begin{proof}
The union of the sets \(\Dx(x_n)\) over \(n\) (with \(k_n\to\infty\)) is the displayed set. A modulus \(\equiv1\ (\mathrm{mod}\ g)\) and \(\ge2\) is at least \(g+1\). The count is at most \(\tau(m)(1+\log_2X)\); the slope is Theorem~\ref{thm:T43}.
\end{proof}

\begin{example}[A natural seed]\label{ex:natural}
The orbit of \(x_0=45828/5=36\cdot1273/5\) has \(\Dx=\{6\}\), \(G=5\), \(\gamma_0=1273\) (Example~\ref{ex:orbit36}). From step 3 on it is in the cycle \(\gamma=1,2,3\) with \(L\) doubling every three steps; the state \(x_3=63504/5\) is a reset state whose class cycle closes (Theorem~\ref{th:F.T51}), so \(G_n=5\) for all \(n\) and \(x_{n+3}=4x_n\) for \(n\ge3\). The orbit of \(20736/5\), by contrast, has \(G=5\) for nine steps and then reaches \(G=1\) (Figure~\ref{fig:growth}).
\end{example}

\section{Universal collapse is false}

\begin{theorem}[Universal \(G\)-collapse is false]\rec{F.T39}\scope{C}\label{thm:T39}
For every \(g\ge5\) with \(\gcd(g,6)=1\) there are collective orbits with \(G(x_n)=g>1\) for every \(n\). Hence the statement ``if \(G_n>1\) then \(G_k<G_n\) for some \(k>n\)'' is false, and the question whether every orbit reaches \(G=1\) is answered in the negative.
\end{theorem}
\begin{proof}
Theorem~\ref{thm:T42} (or Proposition~\ref{prop:T38a}).
\end{proof}

\begin{remark}[Resolution of the collapse problem]\rec{F.L6}\scope{C}\label{rem:L6}
Lemma~\ref{lem:escape} prevents unbounded plateaus and Theorem~\ref{thm:G236} removes \(G=2,3,6\), but neither excludes infinitely many \(G\)-preserving growth events. The collapse question, whether every realised initial value reaches \(G=1\), is therefore settled by Theorem~\ref{thm:T39}: it does not. Random sweeps of natural seeds miss the phenomenon: among \(a<25000\) with \(G=5\) only 10 of 557 persist, and no other \(g\) appears.
\end{remark}

\begin{remark}[Which \(g\) can persist]\rec{F.T41}\scope{C}\label{rem:T41}
Known so far: \(G_\infty\ne2,3,6\) (Theorem~\ref{thm:G236}); no divisor of 24 is a persistent value (Theorem~\ref{thm:G24} below); persistence occurs for every \(g\ge5\) prime to 6 (Theorem~\ref{thm:T42}), for \(g=15,39,51,57,87\) (Theorem~\ref{thm:Pstar}), and for the even value \(g=68\) (Example~\ref{ex:68}). The earlier expectation that persistence requires \(\gcd(g,6)=1\) is superseded by Theorem~\ref{thm:Pstar}. A complete classification of the \(g\) that admit persistent orbits is open (Appendix~\ref{app:open}).
\end{remark}

\section{\texorpdfstring{Exclusion of \(G_\infty\mid24\)}{Exclusion of G(inf) dividing 24}}

\begin{theorem}[\(G_\infty\) cannot divide 24]\rec{F.T45}\scope{C}\label{thm:G24}
If \(g>1\) divides 24, no collective orbit is persistent at \(G=g\). In particular \(G_\infty\notin\{2,3,4,6,8,12,24\}\).
\end{theorem}
\begin{proof}
Suppose the orbit is persistent at \(g\) from \(n_0\).

\emph{\(g=3\).} \(L\) is even, \(G\) collapses iff \(3\mid\gamma\) (Theorem~\ref{thm:gamma}(b),(e)), and \(\gamma\) advances by 1 modulo 3 at every step (Theorem~\ref{thm:gamma}(d)); so \(3\mid\gamma\) within three steps, a contradiction.

\emph{\(g\) even.} Then \(g\mid24\), \(g\ne3\), \(L\) is odd, and \(G\) collapses at the modulus \(M=4g\). Every prime of every \(r\) is prime to \(g\).

\emph{Step 1 (prime classes).} Let \(n>n_0\) and \(q\mid L_n\) a prime. By Theorem~\ref{thm:gamma}(d), \(\gamma\) advances by one modulo 8 if \(g\mid8\), or modulo 24 if \(g\mid24\) with \(3\mid g\); in both cases some \(m\ge n\) has \(g\mid\gamma_m\) (since \(g\mid 8\) or \(g\mid 24\)). By Theorem~\ref{thm:gamma}(e) \(x_m\) is an integer. Because \(q\mid d_0\) for some \(d_0\in\Dx_n\subseteq\Dx_m\), \(q^2\mid a_m\), and since \(b_m=1\), \(q\in\Dx(x_m)\). Persistence gives \(g=G(x_m)\mid q-1\): \(q\equiv1\ (\mathrm{mod}\ g)\).

\emph{Step 2 (growth is large).} Every growth ratio \(r\) at a step after \(n_0\) has all prime factors \(\equiv1\ (\mathrm{mod}\ g)\) by Step 1, so \(r\ge g+1\) and \(r^2\ge(g+1)^2>4g=M\).

\emph{Step 3 (descent).} A plateau has at most \(M-1\) states (otherwise \(M\mid\gamma\), i.e. \(4g\mid\gamma\), forcing \(G=1\)); so growth events are infinite. Let \(a_k\) be the value of \(\gamma\) right after the \(k\)-th growth. The plateau runs \(a_k,\dots,a_k+j\) avoid multiples of \(M\), so \(j\le M-2\) and the next growth carries \(\gamma_{\mathrm{before}}\le a_k+M-2\) to
\[
a_{k+1}=\frac{\gamma_{\mathrm{before}}+1}{r^2}\le\frac{a_k+M-1}{(g+1)^2}<a_k,
\]
the last inequality because \(M-1<a_k\big((g+1)^2-1\big)\) (as \(a_k\ge1\) and \((g+1)^2-1=g^2+2g\ge4g\)). A strictly decreasing infinite sequence of positive integers does not exist.
\end{proof}

The theorem reproves Theorem~\ref{thm:G236}. For \(g=4\) the mechanism is that every prime of \(L\) must be \(\equiv1\ (\mathrm{mod}\ 4)\), forcing \(r\ge5\) and \(25>16\); for \(g=12\), every prime of \(L\) must be \(\equiv1\ (\mathrm{mod}\ 12)\). The factor \(r^2\equiv1\ (\mathrm{mod}\ 8)\) is the reason the modular escape of Lemma~\ref{lem:escape} alone does not settle even \(g\); the descent argument replaces it.

\section{\texorpdfstring{A bounded window for \(\gamma\)}{A bounded window for gamma}}

\begin{theorem}[Persistent orbits live in a bounded \(\gamma\)-window]\rec{F.T47}\scope{C}\label{thm:T47}
Let a persistent orbit have collapse modulus \(M\) (\(M=4g\) if \(g\) is even, \(M=g\) if \(g\) is odd). Then \(\gamma_n\le2M\) for all sufficiently large \(n\).
\end{theorem}
\begin{proof}
(For odd \(g\), \(L\) is even at every state, Theorem~\ref{thm:G-odd}; for even \(g\), \(L\) is odd.) Plateaus have at most \(M-1\) states (else \(\gamma\equiv0\ (\mathrm{mod}\ M)\) and \(G=1\), Theorem~\ref{thm:gamma}(b)), and every growth has \(r\ge2\), so \(a_{k+1}\le(a_k+M-1)/4\). Then \(a_k-\tfrac{M-1}3\le4^{-k}(a_0-\tfrac{M-1}3)\), so \(\limsup a_k\le(M-1)/3\), and \(\gamma_n\le a_k+M-2<2M\) eventually.
\end{proof}

\emph{Meaning.} Up to a bounded factor, a persistent state is the square of the lcm: \(x_n=\gamma_nL_n^2/g\) with \(\gamma_n\) in a finite window. A persistent orbit is a finite-state \(\gamma\)-process driven by arithmetic coincidences; in the regimes of this chapter the process is the cycle \(1,2,3\).

\begin{lemma}[Parity of \(L\) in a persistent orbit]\label{thm:G-odd}
If \(g=G(x)\) is odd then \(L(x)\) is even. If \(g\) is even then \(L(x)\) is odd.
\end{lemma}
\begin{proof}
If \(g\) is odd, some \(d-1\) is odd (else \(2\mid G\)), so some \(d\in\Dx\) is even. If \(g\) is even, every \(d\equiv1\ (\mathrm{mod}\ g)\) is odd.
\end{proof}
